\documentclass[11pt]{article}
\usepackage[margin=0.8in]{geometry}
\usepackage{amsmath,amssymb,amsthm,mathtools}
\usepackage{graphicx}
\usepackage{microtype}
\usepackage[colorlinks=true,linkcolor=blue,citecolor=blue,urlcolor=blue]{hyperref}

\newtheorem{theorem}{Theorem}
\newtheorem{lemma}{Lemma}
\newcommand{\BH}{B(H)}
\newcommand{\KH}{K(H)}
\newcommand{\cI}{\mathcal I}
\newcommand{\cJ}{\mathcal J}
\newcommand{\UU}{\mathcal U}

\title{The two notions of subideal-trace are different}
\author{Automated mathematics run \texttt{fa\_banach\_001}\\Lane 14, model GPT5.6}
\date{9 August 2026}

\begin{document}
\maketitle

\noindent\textbf{Status: candidate full solution, likely valid.}

\section{Source question}

Let $J$ be a two-sided ideal of $\BH$ and let $\cI$ be a $J$-ideal.  Patnaik
and Weiss define
\[
 \UU_J(H)=\{1+A\in\UU(H):A\in J\}
\]
and call a linear functional on $\cI$ a subideal-trace when it is invariant
under conjugation by $\UU_J(H)$.  They later introduce the larger normalizing
class
\[
 \UU^{\cI}(H)=\{U\in\UU(H):UXU^*\in\cI\text{ for every }X\in\cI\}.
\]
On page 9, after Definition~5.11, they ask whether the resulting inclusion of
trace classes can be proper \cite{source}.

\begin{center}
\includegraphics[width=\textwidth]{figures/open_problem_crop.png}
\end{center}

In particular, they ask whether Definitions 5.3 and 5.11 give different
classes of functionals on a subideal which is not a $\BH$-ideal.

\section{Background from the source}

For $S\in J$, the principal linear $J$-ideal has the form
\begin{equation}\label{eq:structure}
 (S)_J=\mathbb C S+JS+SJ+J(S)J.
\end{equation}
When $(S)_J$ is not a $\BH$-ideal, the coefficient of $S$ in this expression
is unique.  Consequently
\begin{equation}\label{eq:tau}
 \tau(\alpha S+X)=\alpha,
 \qquad X\in JS+SJ+J(S)J,
\end{equation}
is a well-defined linear functional, and Example~5.5 of the source proves
that it is $\operatorname{Ad}_{\UU_J(H)}$-invariant \cite{source}.

We also use the Fong--Radjavi criterion quoted on source PDF page 3.  If $T$ is
a compact operator of infinite rank, the principal $\KH$-ideal generated by
$T$ is a $\BH$-ideal if and only if, for some integer $k>1$,
\begin{equation}\label{eq:FR}
 s_{kn}(|T|)=o(s_n(|T|)).
\end{equation}

\section{Proof intuition}

The functional in \eqref{eq:tau} remembers only the scalar coefficient of its
chosen generator.  Conjugation by a unitary of the form $1+J$ changes the
generator only by terms in the remainder of \eqref{eq:structure}, so this
coefficient cannot change.  We choose a generator whose spectrum comes in
opposite pairs.  A block-swap unitary then normalizes the generated subideal
but sends the generator to its negative, forcing the remembered coefficient
to change sign.

\section{Construction and proof}

Let $H_0=\ell^2(\mathbb N)$, let $H=H_0\oplus H_0$, and set
\[
 D=\operatorname{diag}\left(1,\frac12,\frac13,\ldots\right),
 \qquad S=D\oplus(-D),
 \qquad J=\KH,
 \qquad \cI=(S)_J.
\]
Define the block-swap unitary
\[
 W=\begin{pmatrix}0&1\\1&0\end{pmatrix},
 \qquad W(x,y)=(y,x).
\]

\begin{theorem}
The ideal $\cI$ is a $J$-ideal which is not a $\BH$-ideal.  The coefficient
functional $\tau$ from \eqref{eq:tau} is
$\operatorname{Ad}_{\UU_J(H)}$-invariant but is not
$\operatorname{Ad}_{\UU^{\cI}(H)}$-invariant.  Hence the inclusion displayed
after Definition 5.11 of \cite{source} is proper.
\end{theorem}

\begin{proof}
The singular-value sequence of $S$ is
\[
 s(S)=\left(1,1,\frac12,\frac12,\frac13,\frac13,\ldots\right),
 \qquad s_m(S)=\frac1{\lceil m/2\rceil}.
\]
For every fixed integer $k>1$,
\[
 \frac{s_{kn}(S)}{s_n(S)}
 =\frac{\lceil n/2\rceil}{\lceil kn/2\rceil}
 \longrightarrow\frac1k\neq0.
\]
Thus \eqref{eq:FR} fails for every $k>1$.  The Fong--Radjavi criterion shows
that $\cI=(S)_J$ is not a $\BH$-ideal.

Put
\[
 \cI^0=JS+SJ+J(S)J.
\]
The principal-subideal structure theorem gives
$\cI=\mathbb C S\oplus\cI^0$, so the functional
\[
 \tau(\alpha S+X)=\alpha\qquad(X\in\cI^0)
\]
is well-defined and linear.

For completeness, let $V=1+A\in\UU_J(H)$, where $A\in J$.  Then
\[
 VSV^*-S=AS+SA^*+ASA^*\in\cI^0.
\]
Moreover $\cI^0$ is a $J$-ideal, and expansion of
$(1+A)X(1+A^*)$ shows that $V\cI^0V^*\subseteq\cI^0$.  Hence
\[
 V(\alpha S+X)V^*=\alpha S+X'
 \quad\text{for some }X'\in\cI^0,
\]
which proves $\tau(VXV^*)=\tau(X)$ for every $X\in\cI$.  Thus $\tau$ is a
subideal-trace in the sense of Definition 5.3.

On the other hand,
\[
 WSW^*=-S,
 \qquad WJW^*=J.
\]
Conjugation by $W$ therefore maps the $J$-ideal generated by $S$ onto the
$J$-ideal generated by $-S$:
\[
 W\cI W^*=(-S)_J=(S)_J=\cI.
\]
In particular $W\in\UU^{\cI}(H)$.  Nevertheless,
\[
 \tau(WSW^*)=\tau(-S)=-1\neq1=\tau(S).
\]
Therefore $\tau$ is not invariant under $\UU^{\cI}(H)$, proving strictness of
the inclusion.
\end{proof}

Notice also that $W\notin\UU_J(H)$: on the infinite-dimensional antisymmetric
subspace $\{(x,-x):x\in H_0\}$, the operator $W-1$ equals $-2$.  Hence $W-1$
is not compact.  The separating unitary is genuinely available only to the
larger normalizer class.

\section{Verification report}

The proof has four independent checks.
\begin{enumerate}
\item The singular-value ratio tends to $1/k$, so the subideal is not a
$\BH$-ideal.
\item Non-$\BH$-ideality makes the coefficient in \eqref{eq:structure} unique.
\item Conjugation by $1+J$ preserves the coefficient and hence $\tau$.
\item The block swap normalizes the subideal but reverses the coefficient of
$S$.
\end{enumerate}
No numerical experiment is used.  The first two checks use results explicitly
reproduced in the source; the last two are direct computations.

\section{Novelty check and limitations}

The bounded search on 9 August 2026 covered the run indexes, arXiv and web
queries for the exact question and both subideal-trace definitions, and the
locally parsed arXiv corpus.  It found no later explicit answer.  The later
arXiv:2301.09425 cites the survey for subideal terminology but does not discuss
this trace question.

The theorem proves existence and therefore completely answers the source's
properness question.  It does not classify the subideals for which the two
trace classes coincide or compute either full trace space.

\paragraph{Human-review recommendation.}
Check the Fong--Radjavi criterion and the directness of
$\mathbb C S+\cI^0$ first.  Once those source results are accepted, the
normalizer computation is elementary.

\begin{thebibliography}{9}
\bibitem{source}
S. Patnaik and G. Weiss,
\emph{A survey on subideals of operators and an introduction to
subideal-traces}, arXiv:1303.5697 (2013), especially Example~5.5,
Definition~5.11, and page~9,
\url{https://arxiv.org/abs/1303.5697}.
\end{thebibliography}

\end{document}
